\chapter{Where the review says to go next}\label{ch:next}

Chapter~\ref{ch:review} covered what the review found. This chapter covers what it recommends: what the output of a
paper should be, what each kind of prior can learn, which kinds of generalisation are still untested, how to add
uncertainty, and the method it proposes. Nothing here has been built yet. It is the map for after the evaluation.

\section{Which output should come first}

The project has always treated the field at 0.5\um\ as the goal, and the current as a step on the way. The review argues
for the reverse.

\begin{itemize}
\item The 0.5\um\ field and the equivalent sheet current carry the \emph{same} information. Chapter~\ref{ch:magic} showed
that one follows from the other by exact physics.
\item But only the current can be checked in a real experiment: against the lithography mask that made the wires, and
against the current that was driven through them.
\item A 0.5\um\ field on a 0.1\um\ grid can never be measured on a real QDM. Its optical resolution is of the order of the
wavelength of light.
\end{itemize}

\begin{keyidea}
A paper whose ground truth can never exist outside a simulation is hard to defend at a physics venue. Report the
current as the primary result, and the near field as a consequence of it.
\end{keyidea}

\section{What each kind of prior learns}

Chapter~\ref{ch:limits} showed that methods differ only in the prior band. So the question is what each prior knows.

\begin{center}
\small
\begin{tabularx}{\textwidth}{>{\raggedright\arraybackslash}p{3.6cm}X}
\toprule
Prior & What it knows \\
\midrule
Field-domain dictionary (FISTA, CSC-LISTA) & 32 atoms, each a 7.5\um\ patch of field at 0.5\um. Fields of wires are smooth and
overlap, so many atoms are needed everywhere, the codes come out dense, and the prior is weak \\
Current-domain network (MAGIC-UNet) & Currents are compact, piecewise constant, of one amplitude and continuous. That is a
much stronger statement, and it supports filling in detail past $k_c$ \\
\bottomrule
\end{tabularx}
\end{center}

\begin{careful}[The open question]
Is the network recovering structure, or recalling the simulator's shapes? Both test coils are rectangular spirals of the
simulated width. Until it is tested on something outside the simulator's families (curves, diagonals, pads, two current
levels, two layers), there is no evidence either way. Chapter~\ref{ch:ten_results} makes this question sharper, not
smaller.
\end{careful}

One product no other method offers, and it is cheap: for each reconstruction, report which frequencies were supported
by the data (below $k_c$, for the estimated height and noise) and which came from the prior. That is a direct answer to
``how do you handle the ambiguity''.

\section{Which kinds of generalisation are tested}

\begin{center}
\small
\begin{tabularx}{\textwidth}{>{\raggedright\arraybackslash}p{3.4cm}>{\raggedright\arraybackslash}p{2.2cm}X}
\toprule
Axis & Tested? & Status on 16 November \\
\midrule
Standoff, known & Partly & FISTA at 3 and 10\um \\
Standoff, unknown & Partly & Run 2 in 2--4\um; run 3 across 1--10\um; run B across 6--14\um\ on simulation \\
Geometry & No & Two shapes, both in the simulator's families \\
Wire width, thickness & No & \\
Field of view & No & Fixed 140\um\ image \\
Pixel pitch & No & 0.1\um\ input; the network works at 0.2 or 0.7\um\ after averaging \\
Noise level and type & No & One synthetic family; no sensor noise \\
Current amplitude & Implicitly & Each tile is scaled by its own maximum; one current in all data \\
Real measurement & No & \\
\bottomrule
\end{tabularx}
\end{center}

\section{Magnetostatics has no ruler}

This is the most useful idea in the review's recommendations, and you have already seen a rough version of it.

Look at the blur, $e^{-k\,\Delta z}$. If you double every length in the problem, $k$ halves and $\Delta z$ doubles, and the
product $k\,\Delta z$ does not change. So a 2\um\ wire seen from 10\um\ is exactly the same picture as a 1\um\ wire seen
from 5\um, drawn at a different scale. The physics has no built-in length.

\begin{keyidea}[What follows]
Resample every input so that the standoff is a fixed number of \emph{pixels}. Then every standoff becomes the same problem,
with only the wire width in pixels changing. One network trained at that one normalised standoff would serve every
physical standoff, and every pixel pitch too.
\end{keyidea}

Now look back at Chapter~\ref{ch:ten}. Run B's 0.7\um\ pixels did a crude version of exactly this. Run 2 saw a 3\um\
standoff as 15 pixels of 0.2\um. Run B sees a 10\um\ standoff as about 14 pixels of 0.7\um. That is why the same network,
unchanged, works at 10\um\ once the pixels are scaled: in pixels, it is nearly the problem it already knew how to solve.
The full idea needs a height hypothesis to resample with, and the height search below provides one.

\section{Uncertainty, in order of cost}

Nothing in the project outputs an uncertainty yet. From cheapest to most expensive:
\begin{enumerate}
\item \textbf{Seeds.} The spread of scores over several trained copies, and a per-pixel spread of their outputs. Three
seeds exist at 3\um\ (Chapter~\ref{ch:consistency}).
\item \textbf{A height interval.} Plot the data misfit against the hypothesised height, report the range of heights that
fit about equally well, and check on simulated data how often the true height falls inside it.
\item \textbf{The data-supported and prior-filled bands} of the previous section.
\item \textbf{Conformal intervals,} calibrated on simulated data and then checked on COMSOL and real data. Expect them to
break under the gap between simulator and reality; reporting that honestly is useful in itself.
\item \textbf{Posterior samples} from a diffusion prior (candidate B below), for a machine-learning venue.
\end{enumerate}

\section{Physics errors a learned model could absorb or hide}

A network trained on wrong physics does not complain. It quietly learns the error. The review lists the ones to watch:
\begin{itemize}
\item COMSOL's extra decay. A network trained on COMSOL fields would learn it silently. Keeping it out of the simulator
was right, and Chapter~\ref{ch:ten_results} showed what it does to a network that never saw it.
\item The quarter-turn of COMSOL's components. Now traced to the export (Chapter~\ref{ch:comsol}).
\item Mesh noise shared between configurations. Harmless while COMSOL stays out of training; a leak the moment it enters.
\item Sheet currents against conductors of finite thickness, the thickness of the NV layer, and optical blur on a real
QDM. All three act like extra standoff (Chapter~\ref{ch:limits}).
\item MAGIC-UNet's current does not have to conserve charge, because $J_x$ and $J_y$ come from separate decoders.
\item Scaling each tile by its own maximum hides the absolute strength from the network, and makes empty tiles pure
noise. The run 1 bug was a symptom of this.
\item Zero padding, where COMSOL's walls force the normal field to zero.
\end{itemize}

\section{The method the review recommends}

\textbf{Method A: a data-consistent current-domain solver with a search over height.} Five parts, each with a published
ancestor:

\begin{enumerate}
\item \textbf{Unknown: a stream function $g$}, one scalar image, with the current $J=\nabla\times(g\,\hat z)$. Then the
current conserves charge by construction, and $B_x$, $B_y$, $B_z$ at any height follow from $g$ by the Fourier relations
already in \code{simulate_fields.py}.
\item \textbf{Scale normalisation:} resample the input so the hypothesised standoff is a fixed number of pixels (the
section above).
\item \textbf{Network:} a MAGIC-UNet-style encoder and decoder with one output, $g$, used as the learned prior inside a few
unrolled stages. Each stage alternates the closed-form data-consistency step of Chapter~\ref{ch:consistency}, on all
three components, with a learned correction to $g$, in the manner of USRNet. Tiles at least 8 standoffs wide.
\item \textbf{Height:} an outer search. For each candidate height on a grid, resample, reconstruct, then score by the data
misfit and by how much the network had to correct. Report the best height, the interval where the score stays within
a calibrated tolerance, and one reconstruction per supported height when the interval is wide.
\item \textbf{Training:} the simulator with wider variety: curved and diagonal wires, widths over a decade, two current
levels, finite thickness, white and correlated noise, optical blur. COMSOL and real data stay test-only.
\end{enumerate}

\begin{keyidea}[Why this method, for this project]
\begin{itemize}
\item The one experiment we have points to it: one data-consistency step after MAGIC-UNet took 27.0\% to 24.3\%.
\item The run 3 failure and the read-out test of Chapter~\ref{ch:height} say the height must be chosen by
\emph{reconstructing at each candidate height}, not by re-projecting one prediction. Part 4 does exactly that.
\item Every block has a published ancestor (MAGIC-UNet, USRNet, stream-function priors, the joint distance fit of
Setescak et al.), which respects the rule of published parts with a clearly labelled contribution.
\item It reuses the simulator, the Biot--Savart code and the MAGIC-UNet port.
\end{itemize}
\end{keyidea}

The candidates it was compared with:

\begin{center}
\small
\begin{tabularx}{\textwidth}{l>{\raggedright\arraybackslash}X>{\raggedright\arraybackslash}p{3.6cm}}
\toprule
 & Candidate & Verdict \\
\midrule
A & The method above & Build it \\
B & A diffusion prior on $g$, sampled so that it agrees with the measurement, with the height as a variable & Uncertainty
comes naturally, but training is slow and guidance through such a steep blur is unstable. A second paper \\
C & A Fourier neural operator & No clear advantage over a CNN after resampling \\
D & An untrained network with the exact physics as its loss, height as a free parameter (Dubois et al.) & No simulator gap
at all. Build it anyway as a baseline \\
E & Classical total variation on the stream function, at a grid of heights & A baseline and the identifiability test in
one. Build it anyway \\
\bottomrule
\end{tabularx}
\end{center}

\section{The one claim a paper could make}

\begin{keyidea}[The claim, as the review writes it]
The standoff of a magnetic image can be inferred, with an honest interval, from the image itself once a learned current
prior and exact field physics are combined. Doing so gives better current and near-field reconstructions than methods
that are told the standoff, across a tenfold range of standoffs, and the reconstruction states which part of it the data
support.
\end{keyidea}

What a reviewer would demand, and what exists: baselines (Fourier inversion, total variation on the stream function, an
untrained network, MAGIC-UNet as published, FISTA; we have the last two); data sets (held-out simulated families, ten or
more new finite-element geometries with correct physics, one real QDM measurement with a calibrated standoff; we have two
shapes with an anomalous decay); ablations of every ingredient; three or more seeds; and a test protocol written before
the final runs. The review's realistic venues are Physical Review Applied or IEEE Transactions on Computational Imaging,
with a real measurement.

\section{The five biggest risks}

\begin{enumerate}
\item \textbf{The benchmark is not valid evidence.} Two shapes, a non-physical decay, inconsistent components, a noisy
target. Fix the data before comparing methods further; a real measurement for the paper.
\item \textbf{The learned prior has memorised the simulator's shapes and widths.} Every gain so far is on coils the
simulator was built around. Test on held-out families and new shapes, and keep the untrained-network and total-variation
baselines next to every result, since they have nothing to memorise.
\item \textbf{The height is not identifiable tightly enough to be useful.} The valley is shallow at 3\um\ and will be
shallower higher up, and width and thickness bias it. Report an effective standoff with an interval. Keep the fallback
paper in view: reconstruction given a calibrated standoff, with its sensitivity to errors in that standoff.
\item \textbf{The novelty is thinner than it looks.} MAGIC-UNet already maps field to current, possibly over a standoff
range; Setescak et al.\ already fit the sensor distance jointly. Re-read both in full.
\item \textbf{Statistical fragility.} One seed, many variants on 15 pairs, choices informed by earlier scores. Write the final
protocol first, freeze a new test set no method has seen, use three seeds, and report every run.
\end{enumerate}

A sixth, practical one: one GPU and one person. Method A fits that budget; a diffusion prior and a wide architecture
search do not.

\begin{defend}[If your professor asks: what would you do with three more months?]
Build method A: a stream-function output so the current conserves charge, inputs resampled so the standoff is a fixed
number of pixels, a learned prior inside a few data-consistency steps, and a search over height that reconstructs at
each candidate and reports an interval. Test it on new geometries and on one real QDM measurement with a calibrated
standoff, against an untrained network and total variation, which have no simulator to memorise.
\end{defend}

\begin{defend}[If your professor asks: what is genuinely new here?]
Not learning field to current, not using all three components, and not fitting the distance jointly as an idea: all of
those exist. What is open is inferring the standoff for the whole image with a learned current prior and exact physics,
the one-sided analysis of the ambiguity with a measured height resolution, and a reconstruction that says which
frequencies came from the data and which from the prior.
\end{defend}
